\documentclass{article}
\usepackage[UTF8]{ctex}
\usepackage{amsmath,amssymb,amsfonts}
\usepackage{enumitem}
\usepackage{geometry}
\geometry{a4paper, margin=1in}
% === 错题本排版优化（不参与题目标识） ===
\setlength{\parindent}{0pt}
\setlength{\parskip}{0.35em plus 0.10em minus 0.08em}
\setlength{\emergencystretch}{3em}
\setlist{itemsep=0.25em, topsep=0.40em, parsep=0pt}
\allowdisplaybreaks[2]
% === 错题本排版优化结束 ===
\begin{document}
\section*{天津大学2026年高等代数真题}
\section*{53．天津大学2026年高等代数试卷}

\subsection{题目1}
1．称 \(f \in \mathbb{R}\left[x_{1}, x_{2}, \cdots, x_{n}\right]\) 为反对称多项式，如果对 \([1,2, \cdots, n]\) 的任意置换 \(\sigma\) ，均有

\[
f\left(x_{\sigma(1)}, x_{\sigma(2)}, \cdots, x_{\sigma(n)}\right)=\varepsilon_{\sigma} f\left(x_{1}, x_{2}, \cdots, x_{n}\right) .
\]

其中 \(\varepsilon_{\sigma}\) 表示 \(\sigma\) 的符号．令 \(f\left(x_{1}, x_{2}, \cdots, x_{n}\right)\) 为反对称多项式．证明：

\begin{itemize}
  \item[（1）] 对任意的 \(i \neq j\) ，有 \(\left(x_{i}-x_{j}\right) \mid f\left(x_{1}, x_{2}, \cdots, x_{n}\right)\) ．
  \item[（2）] \(\Delta=\prod_{1 \leq i<j \leq n}\left(x_{i}-x_{j}\right)\) 为反对称多项式．
  \item[（3）] 存在对称多项式 \(g\) ，使得 \(f=g \Delta\) ．
\end{itemize}

\subsection{题目2}
2．已知 \(A, B, C, D\) 均为 \(n\) 阶方阵，证明：

\begin{itemize}
  \item[（1）] \(\left|\begin{array}{ll}A & B \\ B & A\end{array}\right|=|A+B||A-B|\) 。
  \item[（2）] \(\left|\begin{array}{llll}A & B & C & D \\ B & A & D & C \\ C & D & A & B \\ D & C & B & A\end{array}\right|=|A+B+C+D||A+B-C-D||A-B+C-D||A-B-C+D|\) ．
\end{itemize}

\subsection{题目3}
3．设 \(\mathrm{Q}[A]=\{f(A) \mid f(x) \in \mathrm{Q}[x]\}\) ，其中 \(A=\left(\begin{array}{ll}0 & \sqrt{2} \\ 1 & 0\end{array}\right)\) ，证明： \(\mathbb{Q}[A]\) 中的非零矩阵均为可逆矩阵。

\subsection{题目4}
4．设 \(V\) 是 \(n\) 维线性空间，\(e_{1}, e_{2}, \cdots, e_{n}\) 是 \(n\) 个向量，若 \(V\) 中任一向量均可由 \(e_{1}, e_{2}, \cdots, e_{n}\) 线性表出， i明：\(e_{1}, e_{2}, \cdots, e_{\mathrm{n}}\) 是 \(V\) 的一组基．

\subsection{题目5}
5．化下列实二次型为标准形：\(f\left(x_{1}, x_{2}, \cdots, x_{\mathrm{n}}\right)=n \sum_{i=1}^{n} x_{i}^{2}-\left(\sum_{i=1}^{n} x_{i}\right)^{2}\) 。

\subsection{题目6}
6．设 \(V\) 是 \(n\) 维线性空间，\(V^{*}\) 是 \(V\) 的对偶空间，若 \(W \subset V^{*}\) 是 \(r\) 维子空间，证明：

\[
W^{\perp}=\{v \in V \mid l(v)=0, \forall l \in W\} \subset V
\]

是 \(n-r\) 维线性子空间。

\subsection{题目7}
7．设 \(T: V \rightarrow V\) 是实数域 \(n\) 维线性空间 \(V\) 上的一个线性算子，且 \(T^{3}=T\) ，其中 \(n>1\) 。\\
（1）证明：存在 \(V\) 的一组基 \(\alpha_{1}, \alpha_{2}, \cdots, \alpha_{n}\) ，使得 \(T\) 在这组基下的矩阵为对角矩阵。\\
（2）如果存在一个 \(n\) 阶实方阵 \(A\) ，使得 \(T\) 在 \(V\) 的任意一组基下的矩阵均为 \(A\) ，证明：\(A\) 为 \(E_{n},-E_{n}\)或零矩阵。

\subsection{题目8}
8．设 \(n\) 阶复矩阵 \(A\) 有 \(n\) 个不同的特征值，其中 \(n>1\) 。\\
（1）设 \(B\) 为一个 \(n\) 阶复方阵，且 \(B\) 与 \(A\) 有完全相同的特征多项式。证明：存在两个复方阵 \(P, Q\) ，使得 \(A=P Q\) 且 \(B=Q P\) 。\\
（2）\(C\) 为 \(n\) 阶复方阵，且 \(A C=C A, C^{n}=O\) ，求矩阵 \(C\) 。

\subsection{题目9}
9．设 \(T: V \rightarrow V\) 是欧氏空间上的一个正规算子。

\begin{itemize}
  \item[（1）] 证明：\(T\) 的伴随算子的核空间与 \(T\) 的核空间相等．
  \item[（2）] 若 \(T^{2}=T\) ，则 \(T\) 必为对称算子。
\end{itemize}
\end{document}
